\cotefile{156-3}
\scanpage{34}\mspage{33}
% L1–3
\deleted{\unsure{7|8}}\unsure{8}) $\forall\,X\in M$, soit $X'$ son ``complémentaire'' \deleted{\unsure{8|6}} Comme $\{X,X'\}$ \unsure{modéré}, alors $\{\overline X,\overline{X'}\}$ \deleted{$\mathrm{Fr}\,X$ est \unsure{fini} $\unsure{\{\overline X,\overline{X'}\}|\mathrm{Fr}\,X}\in M$} \inserted{aussi, on peut} former
\[
\left\{\begin{aligned}
\mathrm{Fr}\,X &\overset{\text{def}}{=} \overline X\wedge\overline{X'}\\
(\mathrm{Fr}\,X)' &= X^{\circ}\vee X^{\prime\circ}
\end{aligned}\right.
\]
% slanted
\margin{propriété spéciale au cas envisagé du modèle de \unsure{formes} \uline{de} \uline{dim $1$}}
% L4–5
Ceci dit $\{X,\deleted{\unsure{X'}}\unsure{Y|Y'}\}$ modéré \deleted{\unsure{ssi}} $\{\mathrm{Fr}\,X,\mathrm{Fr}\,Y\}$ est modéré $\Longrightarrow$ $\mathrm{Fr}\,X\cup\mathrm{Fr}\,X$ est modéré.

% L6–10
\deleted{\uline{NB}} \deleted{Soit $X$ \unsure{modéré}, alors} \deleted{\inserted{On dit que $X$ \unsure{($\in M$)|(de $M$)} \unsure{est} \uline{rare} si}} \deleted{\unsure{si} $X^\circ=\varnothing$ i.e. $\overline{X'}=1$ ($X'$ \uline{dense}). Dans le cas d'espèce, $X$ est \unsure{rare|un} \unsure{sss|ssi} $X$ est fini (en \unsure{tant} que partie de $L$), ou encore sss $\mathrm{Fr}\,X$}

% L11–12
Dans le cas d'espèce, conditions équivalentes sur $X\in M$
% L13 (a)
\[
\begin{aligned}
\text{a)}\quad & X\subset\deleted{\unsure{\mathrm{Fr}(X)}}X^{\bullet}
&&\left(\begin{aligned}&\text{i.e.}\quad X^{\circ}=\varnothing\quad\unsure{\&}\ \text{i.e.}\quad\overline{X'}=1\\[-2pt]&\hspace{8em}\text{i.e. }X'\text{ ``dense''}\end{aligned}\right)\\
% L14 (b)
\text{b)}\quad & X=X^{\bullet}
&&\left(\begin{aligned}&\text{i.e. }X\text{ fermé \unsure{rare} \quad i.e. }X'\text{ ouvert}\\[-2pt]&\hspace{12em}\text{dense}\end{aligned}\right)
% L15 (c)
\end{aligned}
\]
c) $X$ est fini en \unsure{tant} que partie de $L$

\begin{center}\rule{1.6cm}{0.45pt}\end{center}

% L16–20
Considérons un \unsure{ens} ordonné $M$ satisfaisant les conditions 1) à 7) (avec 4) \unsure{renforcé} \unsure{marginalement} comme il a été dit~; dans la notation algébrique, on a
% L21–24
\diagram{\begin{tikzcd}[ampersand replacement=\&,row sep=0.1ex,column sep=1.8em,
  cells={nodes={anchor=base west}},execute at end picture={
  \draw ([xshift=-0.9em]alJoin.north west) -- ([xshift=-0.9em]alOrder.south west);
  \draw[->] ([xshift=0.1em]alOrder.north west)
      .. controls ([xshift=-0.45em]alMeet.west) and ([xshift=-0.7em]alJoin.west)
      .. ([xshift=-0.45em]alJoin.north west)
      .. controls (alJoin.north west) and ([xshift=0.5em]alJoin.north west)
      .. ([xshift=1.1em,yshift=0.25em]alJoin.north west);
  \draw ([xshift=2.6em]alOrder.north west) to[bend right=8] (alOrder.north west);
  \draw (alOrder.south west) to[bend right=8] ([xshift=2.6em]alOrder.south west);
}]
X+Y=(X\vee Y)\wedge(X\wedge Y)' \& |[alias=alJoin]| X\vee Y=X+Y-XY\deleted{\unsure{Y|X}}\\
X\cdot Y=X\wedge Y \& |[alias=alMeet]| X\wedge Y=X\cdot Y\\
0=\text{plus petit él.} \& |[alias=alOrder]| X\leq Y\Longleftrightarrow X\cdot Y=X\\
1=\text{plus gr\textsuperscript{d} él} \&
\end{tikzcd}}
\scanpage{35}\mspage{34}
% L1–2
Soit $\deleted{X}A\in M$. On va ``relativiser'' \unsure{tout} p.r. à $\deleted{X}A$. Soit \deleted{$M_X=M_{\leq X}$,}
% L3
\[
M_{\deleted{X}A}=\deleted{\unsure{\{}}M_{\leq\deleted{X}A}=\{Y\in M\mid Y\subset\deleted{X}A\}
\]
% L4–5
$M_{\deleted{X}A}$ muni de la relation d'ordre induite par $M$

% L6
$\{X,\deleted{\unsure{Y}}Z\}$ est mod dans $M_{\deleted{X}A}$ \unsure{sss|ssi} il l'est dans $M$

% L7–9
\[
\begin{aligned}
X\in M_{\deleted{\unsure{X}}A}\text{ est fermé}\quad &\text{sss}\qquad X=\overline X\wedge\deleted{X}A\\
&\text{i.e. sss}\quad\exists\,\deleted{\unsure{Y}}Z\in M_f,\ \{\deleted{\unsure{Y}}Z,\deleted{X}A\}\ \text{mod et}\\
&\hspace{10em}X=\deleted{\unsure{Y}}Z\wedge\deleted{X}A
\end{aligned}
\]

% L10–11
Je dis que les conditions 1) à 7) sont encore satisfaites~:
% L12, 1)
\[
\text{1)}\quad\varnothing_{M_{\deleted{\unsure{X}}A}}=\varnothing_{M_{\deleted{\unsure{\leq X|X|A}}}},\qquad 1_{M_{\deleted{\unsure{X}}A}}=X
\]
% L13, 2)
2) \deleted{$\unsure{X\wedge Z}$} trivial

% L14–15, 3)
3) item car \deleted{$\unsure{X\overset{M}{\wedge}Z}$} $X\overset{M}{\wedge}Y$ et $X\overset{M}{\vee}Y$ sont dans $M_A$ donc ce sont des bornes inf et sup dans $M_A$.

% L16–23, 4)
4) \deleted{$X\unsure{<|\subset}A$,} $X\in M_A$, considérons les $Y\in M_{\unsure{Y|y}}$ en position mod. p.r. à $X$, telles que $Y\wedge X=\varnothing$ i.e. \deleted{les $Y\subset\unsure{X'|X}$} on a $X\subset A$, on considère les $Y\subset A$ tels que $\{X,Y\}$ mod. et $X\wedge Y=\varnothing$. \unsure{Ces|Les} $Y$ sont $\leq X'$ (déf de $X'$) et $\leq A$, donc $\leq X'\wedge A$, qui a les propriétés voulues, et c'est \unsure{la|le} plus \unsure{grande|grand} de tous, OK

\scanpage{36}\mspage{35}
% L1–4
Considérons des $X_i\in M_A$ \unsure{mutuellement} modérés, et les opérations $\unsure{\complement|C}_A,\wedge,\vee$ \deleted{i.e.} \unsure{i.e|(i.e} des $X_i\subset A$ mut. mod (et mod p.r. à $A$ \unsure{autom.}) on regarde les op.
% L5–7
\deleted{$\complement_A\ \wedge,\vee$ \unsure{et} (binaires)\unsure{,} \unsure{$\overline{\phantom X}$}} \deleted{\inserted{binaires \unsure{, $\overline{\phantom X}$}}} $\complement_{M_A}$ (unaire) $\wedge,\vee$ (binaires) et on adjoint $0_{\deleted{\unsure{A}}M_A},1_{\deleted{\unsure{A}}M_A}$, on voit que c'est pareil que
% L8–9
\deleted{\unsure{de l'adjoindre|d'y adjoindre}} de prendre $\complement_M,\wedge,\vee,0_M,1_M\quad\overline{\phantom X}$ et de remultiplier par $A$. \deleted{\unsure{D}}Mais
% L10–12
alors on voit que l'on trouve un idéal d'une alg. de Boole, c'est une alg. de Boole \ldots.

% L13–14
\deleted{\unsure{7|5}} Il faut prouver que
\diagram{\begin{tikzcd}[ampersand replacement=\&,row sep=1.4em,column sep=1em]
\& \underbrace{\{X,Y\}\ \text{mod}}\arrow[d,Rightarrow]\\
\deleted{\unsure{\Longrightarrow}}\ \{\complement_{M_A}X,\complement_{M_A}Y\}\ \text{mod,}\quad\text{mais} \& \{\underbrace{\complement_M X}_{X},\underbrace{\complement_{\unsure{M|m}}Y}_{Y}\}
\end{tikzcd}}
% L15–17
mod, donc comme \deleted{$A$ \unsure{s.}} les deux \unsure{termes} sont mod. avec $A$) $\{A\wedge X',A\wedge Y'\}$ modéré, OK.

% L18
5) Trivial

% L19–22
6) Soit $X\subset A$, les fermés dans $A$ sont les $Z=\deleted{\unsure{X|Z}}Y\wedge A$, avec $Y\in M_f$, $\{Y,A\}$ mod \deleted{\unsure{les} $\unsure{Z}$} \inserted{les $Z$} qui \unsure{contient|contiennent} $X$ \deleted{sont \unsure{ceux|les $Y$} qui \unsure{contiennent}} \inserted{\unsure{correspondent} \deleted{\unsure{à|i.e.} $Y\supset X$}} avec $Y$ t.q. $Y\supset X$ i.e. $Y\supset\overline X$, et le
\scanpage{37}\mspage{36}
% L1–2
plus petit possible est $Y=\overline X$, d'où \deleted{$Z=\overline X\unsure{,|\wedge}$}
\[
\overline X^{(M_A)}=\overline X^M\wedge A
\]
% L3–6
Soient $X,Y\in M_{\deleted{\unsure{X}}A}$, $\{X,Y\}$ modéré\unsure{)|,} alors
\[
\begin{aligned}
\overline{X\deleted{\unsure{\wedge}}\vee Y}^{M_A}
   &=\underbrace{\overline{X\vee Y}^M\wedge A}\\
   &\hspace{-.4em}(\overline X\vee\overline Y)\wedge A
       =\underbrace{(\overline X\wedge A)}\vee(\overline Y\wedge A)\\
   &\hspace{12em}\overline X^A\vee\overline Y^A
\end{aligned}
\]

% L7–11
7) $X,Y\in M_{\inserted{A}}^{\deleted{\unsure{A}}}$, $\{X,Y\}$ modéré) \deleted{\unsure{alors|donc}} \deleted{donc $\unsure{\{\overline X^{M},\overline Y^{M}\}}$ \unsure{aussi}} $\{X,Y,A\}$ mod. donc $\{\overline X,Y,A\}$ mod donc $\{\overline X\wedge A=\overline X^A,Y\unsure{\{|\}}$ modéré.

\uline{OK}

% L12–15
\uline{Composantes}\qquad On dit qu'un $M$ est connexe si \deleted{$\forall\,X$ fermé et ouvert,} $\forall\,X\in M$, $X$ fermé et $X$ ouvert
\[
\Longrightarrow\quad X=\varnothing\quad\text{ou}\quad X=\unsure{1|\underline{1}}
\]

% L16–21
\deleted{\unsure{N|S}}\uline{NB}\qquad Soit $X$ tel que $X$ fermé et ouvert, On s'attend qu'il soit mod. avec \deleted{\uline{\unsure{$\forall A$|$A$}}} $Y$. S'il en est ainsi, Alors les $X$ en question engendrent une algèbre de Boole. \deleted{Si elle est finie, \unsure{les|on} \unsure{peut|points}} \deleted{\unsure{du spectre}} \inserted{les éléments de} son spectre jouent le rôle de comp. conn. Si l'alg.
\scanpage{38}\mspage{37}
% L1–4
de Boole en question est finie, alors les comp. connexes sont les \deleted{\unsure{$Y$|$X$}} $X\ne\varnothing$ \uline{minimaux} qui soient ouverts et fermés. [en supposant $0\ne1$
% L5–7
i.e. $M$ non réduit à un p\textsuperscript{t} --- sinon on a \deleted{une} \unsure{mor\textsuperscript{t}|moral\textsuperscript{t}|ess\textsuperscript{t}} une structure ``vide'' \ldots\ ]

% L8–10
\uline{Ex}\qquad Reprenons $M=\mathrm{Mod}(\unsure{L|\mathcal L})$, et $M_A$ où $A\in M=\mathrm{Mod}(L)$, donc $A$ partie modérée d'un tronçon.

% L11–12
On aura donc
\[
\begin{aligned}
A=X_{\unsure{\varphi|\Phi}} &,\quad
\varphi\in{\unsure{\mathrm{Drop}|\mathrm{Drp}}}^{\unsure{*|K}}
  \bigl(\unsure{\mathrm{Tro}|\mathrm{Tr}_0}(L)\bigr)\\[-0.6ex]
 &\mkern-8mu\inserted{=\bigcup_i T_i}
\end{aligned}
\]
% L13–14
Je dis que les comp. connexes de $A$ sont les \unsure{ps.tronçons|pr.tronçons} composant $T_i$.
% L15–17
En effet, on a bien \unsure{sur|sûr}
\[
A=\bigvee_i T_i\quad\text{dans }M,\quad\deleted{\unsure{\&}}\quad\text{et}\quad T_i\wedge T_j=\varnothing
\]
% L18–22
si $i\ne j$, il suffit donc de prouver que a) Chaque $T_i$ est ouvert et fermé dans $A$, et b) Chaque $T_i$ est connexe, i.e. ne peut se décomposer en deux parties qui \deleted{\unsure{soit|soient}} disjointes et \inserted{rel\textsuperscript{t}} fermées $\ne\varnothing$.

\scanpage{39}\mspage{38}
% L1–4
Pour a), il suffit de prouver que $T_i$ est fermé dans \deleted{$\unsure{T|X}$} $A$ (les $T_j$, $j\ne i$, l'étant aussi, $T_i$ est aussi ouvert \ldots\ ) \unsure{Si} \deleted{\unsure{$T_i$ est f}} i.e. que
\[
\underbrace{\overline{T_i}}_{\inserted{\substack{\text{``adh'' au}\\\text{sens modéré}\\\text{bien sûr!}}}}\,\unsure{\cap|\wedge}\, A=\unsure{T|\Pi}_{i},
\]
% L5–6
C'est immédiat

% L7–12
\uline{NB}\qquad Si $T=T_i$ est un \unsure{ps.tr.|pr.tr.}, son ``adh.'' modérée n'est autre que le ps.tronçon \inserted{fermé} associé ($=T$ si $T$ réduit à un p\textsuperscript{t}) Il suffit de voir ce qui se passe quand $T$ n'est pas fermé, en un \deleted{$a\in$} \deleted{$\unsure{\partial T}$} \deleted{$\mathrm{Fr}(T)=\unsure{\dot T|T^\bullet}$ tel que}
\[
a\in\partial T\setminus T\quad\ldots.
\]

% L13–14, b)
b) $T$ \uline{connexe}\quad \inserted{Si $T=\{\unsure{a|0}\}$ trivial. Sinon $T$ un \deleted{\unsure{ps.}} tronçon.} \uline{NB}\quad $(M_{\unsure{L|\mathcal L}})_T=M_T$,
% L15–17
\unsure{en regardant|on regarde} $T$ comme un tronçon en lui-$\overline{\text{m}}$. \deleted{\unsure{Alors}} Donc \unsure{OPS} $T=L$. Mais Si $X\subset\deleted{\unsure{L|T}}\unsure{T|L}$ modéré,
% L18–19
alors \unsure{:} \deleted{$\unsure{X}$} $X=X_{\unsure{\varphi|\Phi}}=\bigcup_{i\in\unsure{I|F}}T_i$, alors \deleted{\unsure{Fr}} la frontière modérée est la réunion des \deleted{$\unsure{\partial}$} $T_i^{\unsure{\bullet|\prime}}$,
% L20–22
elle est vide \unsure{sss|ssi} les $T_i^\bullet$ sont vides. $\unsure{\checkmark}$Mais $T_i^{\unsure{\bullet|\prime}}=\deleted{\text{vide}}\varnothing\Longrightarrow T_i=L$ donc $X=\varnothing$ ($I=\varnothing$) ou $X=L$ ($I\ne\varnothing$) OK.

\scanpage{40}\mspage{39}
% L1–3
\uline{NB}\qquad Soit $(X_\alpha)$ famille \inserted{finie} de parties modérées de $L$, \deleted{\unsure{mut\textsuperscript{t}.}} \unsure{mutuell\textsuperscript{t}} modérées i.e telles que
% L4
\[
S=\bigcup_\alpha X_\alpha^{\unsure{\bullet|\circ}}
\]
% L5–6
soit une partie \unsure{tot} ordonnée~:
\[
t_1<t_2<\ldots<t_n
\]
% L7–12
Alors l'Alg. de Boole engendrée par les $X_\alpha$ \deleted{est \unsure{consti|enti}} admet comme \deleted{composantes} \deleted{\unsure{éléments} du spectre (quand on les interprète comme parties de $L$,} ``composants'' les \unsure{ens} modérés \inserted{(et \unsure{ps.tronçons|pr.tronçons})} suivants (rangés par ordre croissant pour $<$)
% L13–14
\[
\scriptstyle
\underset{\substack{\vert\\\inserted{\text{si }t_1\notin\partial L}}}{[L_{<t_1}]}\unsure{,|a}\;
\{t_1\},\ ]t_1,t_2[,\ \{t_2\},\ ]t_{\unsure{1|2}},t_3[,\ \{t_3\},\ldots,\
\deleted{\unsure{\{t\}}}\ \{t_n\},\
\underset{\substack{\rule{3em}{0.35pt}\\\inserted{\text{si }t_n\notin\partial L}}}{[L_{>t_n}]}
\]
% L15–18
Ainsi, cette alg. de Boole est connue quand on connait $S$, i.e. les points $t_1<\ldots<t_n$ (\unsure{ou la|ou} ``décomposition'' \inserted{ou ``subdivision''} du tronçon \unsure{$L$|$\Lambda$|$A$}).

\scanpage{41}\mspage{40}
% L1–3
\uline{Caractérisation d'un} \uline{intervalle} \deleted{\unsure{réseau-interval|réseau-intervalle}} \uline{ordonné}, \uline{en \deleted{\unsure{tant}} que réseau}

% L4–5
\uline{par la donnée} de
\diagram{\begin{tikzcd}[ampersand replacement=\&,row sep=0.75em,column sep=0.15em]
S \& {}\subset\unsure{\mathrm P|\mathcal P}_2(\unsure{L|\mathcal L})\\
\makebox[0pt][c]{\(\unsure{\mathrm{Drop}|\mathrm{Drp}}_2(L)\)}\arrow[u,equal] \&
\end{tikzcd}}
% L6–9
et pour $\unsure{\varepsilon|\epsilon}=\{a,b\}\in S$, la décomposition de \deleted{$L$ \unsure{en}} $\unsure{\complement|C}_{\mathrm{mod}}\,\varepsilon=\{x\unsure{\in|+}L\mid\{\unsure{a;x|a',x|a,x}\},\{b,x\}\in S\}$ en ``composantes connexes'' (par relation $R_\varepsilon$)

% L10–18
J'ai commencé à faire l'exercice, mais j'ai l'impression que ça n'en vaut pas la peine de l'expliciter jusqu'au bout. Ça doit avoir été fait depuis longtemps (Hilbert~?) dans le cas totalement ordonné \inserted{(et divisible)} (en termes \unsure{tout au|tant au|tout aux} moins de la relation~: $\unsure{s|\varepsilon|c}$ est entre $a$ et $b$ \ldots\ )
